\section{What justifies an inquiry-relative classification?}\label{sec:assessment}

An inquiry gives an investigator a reason to distinguish cases. Evidence then
determines whether a proposed grouping supports that distinction. For an
expressive modifier such as \mention{damn}, the grammatical question concerns
its place inside the noun phrase; the semantic question concerns its
contribution to interpretation. A useful answer to either can leave the other
open. Calling both answers classifications doesn't give them equal evidence,
but it gives them the same obligation to state what they accomplish.

\subsection{Bearers, grouping, contribution, and alternatives}

Four questions organize that obligation. What bears membership? What groups
the members? What does membership help an investigator distinguish, infer, or
explain? And what alternative would serve the same inquiry? Each proposal
also needs a scope: a language or community, relevant conditions, and, where
the claim concerns change, a period. These questions apply to the three
worked classifications and to every proposed extension of the family.

The \term{bearer} is the object being classified. In the worked cases it is
a conventionalized sign-type, an expression in an identifiable conventional
use. A phonologist may instead classify its realizations, and a historian
may classify the changes linking it to earlier signs. A statement about the
frequency of realizations doesn't immediately become a statement about the
number of types or the rate of change. Specifying the bearer prevents those
transfers before they enter an argument.

The \term{grouping basis} is the evidence or relation used to assign members.
Grammatical distribution, conventional expressive contribution, common
ancestry, and participation in a socially recognized repertoire are different
possibilities. A grouping can be categorical or graded; it can be supported
by one diagnostic or several. Its shape should follow the evidence relevant
to the inquiry. A demand for a unique mechanism or a sharp boundary would
exclude proposals before their intended contribution had been assessed.

The contribution concerns what investigators can do with the grouping. A
descriptive classification may organize recurrent contrasts across cases;
a projective one supports expectations about cases or properties beyond
those used to assign membership; an explanatory one connects differences
to independently supported processes. These achievements have different
requirements. A label that restates its admission criterion has identified
a set, but its further scientific value remains to be shown.

An \term{alternative} can be a broader category, a narrower subtype, a
continuous dimension, or a different organization of the same evidence.
For expressive \mention{damn}, an undifferentiated interjection category
competes with a classification that groups it with other expressives despite
their grammatical differences. For recipient responses, a broad interactional
class competes with separate action subtypes. The appropriate comparison
depends on the question; a smaller vocabulary isn't invariably a better
account.

Table~\ref{tab:inquiry-record} shows how the questions constrain a proposed
classification. The final column identifies what would still require
evidence. The example concerns expressive contribution rather than a claim
that all signs in the traditional interjection domain share it.

\begin{table}[htbp]
\centering
\small
\caption{Assessing the expressive classification}
\label{tab:inquiry-record}
\begin{tabular}{@{}>{\raggedright\arraybackslash}p{.17\textwidth}
                  >{\raggedright\arraybackslash}p{.36\textwidth}
                  >{\raggedright\arraybackslash}p{.39\textwidth}@{}}
\toprule
Question & Proposed answer & Evidential consequence \\
\midrule
Bearer & Conventional expressive uses, including standalone and prenominal
\mention{damn} & Establish the uses independently of the particular target
being assessed. \\
\addlinespace
Grouping & A conventional contribution distinguishable from descriptive
content & Host-proposition preservation alone is insufficient; descriptive
supplements can preserve it too. \\
\addlinespace
Contribution & Preserve a semantic generalization across different
grammatical constructions; investigate further expressive diagnostics &
Distinguish an attested contrast from a prediction about an unexamined sign
or environment. \\
\addlinespace
Alternative & An undifferentiated interjection class, item-specific
descriptions, or a broader expressive classification & Compare their
coverage and consequences for the semantic question, with comparable
evidence. \\
\bottomrule
\end{tabular}
\end{table}

\subsection{From candidate to supported classification}

A candidate earns consideration by making those questions answerable. A
supported classification has evidence for the contribution it claims.
Comparative merit requires a further judgment about alternatives. The
existing grammatical, semantic, and interactional traditions supply
substantial descriptive resources, but the three reconstructions in this
paper inherit their evidence selectively. Their shared notation supplies
no evidence of its own.

The distinction matters when extending the family. A phonological stratum
with a clear bearer and plausible generalizations deserves investigation
before those generalizations have been demonstrated. Its status at that
point is prospective. The same is true of an untested projection from
\intj{syn} or \intj{prag}. Equal standards of assessment permit unequal
evidential support; they require that the difference in support be made
visible.

Even a well-supported classification has a limited reach. Evidence for
the expressive contribution of prenominal \mention{damn} doesn't settle
the contribution of repair \mention{huh?}. An analysis of one sequential
action doesn't establish a common next-turn prediction for all interactional
resources. Generalization requires evidence for the grouping at the breadth
claimed. Where only a subtype is informative, the larger category may
remain useful for description without acquiring that subtype's predictions.

\subsection{Projection and identity}

\term{Projectibility} concerns warranted expectations beyond the cases in
hand \citep{Goodman1955}. A projective claim specifies what has been observed,
what is expected, and over which bearers and conditions. Its \term{warrant}
is the evidence supporting that expectation. A precise specification can
have little warrant; a familiar association can be useful even while its
cause remains unknown. These distinctions separate the content of a claim
from its evidential and explanatory status
\citep[11--12]{reynolds2026projectibilityHistory}.

Suppose independently assessed source contexts identify an expressive
use, and the target is its behaviour under repetition in another context.
Repetition can't also have been the case-level evidence used to admit that
use to the source class. Potts's expressive diagnostics provide possible
targets of this kind \citep[166--167]{potts2007}. Their availability in a
theory doesn't establish their success for every proposed member.

The same separation constrains sign-type identity. Standalone and
prenominal \mention{damn} can be distinguished through recurrent
constructional and interpretive contrasts. But a target failure can't
automatically create a new sign-type. If determiner resistance is the
prediction, encountering a determiner doesn't by itself license excluding
the case as conversion. The identification rule has to be fixed from
source evidence; an independently supported change in that rule creates
a revised claim requiring fresh assessment.

An informative restriction also needs relevant opportunities. Failure
to observe an inflected form says little if the sampled environments
wouldn't elicit inflection under a competing analysis. For a grammatical
projection, the design should name the source construction, the target
construction, and the competing expectation. For an interactional one,
it should identify the action and the environments in which recipients'
treatment would count as support or counterevidence. Such specifications
are needed for a test; the classifications themselves don't supply them.

\subsection{Representation and explanation}

A category can organize evidence without adding a new fact to a complete
description of its members. If its label is computed deterministically from
fully observed features, adding that label to those same features adds no
information. The useful comparison concerns an investigator who has limited,
noisy, or incomplete observations. Grouping can support pooling across cases,
reduce the number of separate assumptions, or represent dependencies that
an alternative overlooks. Those benefits are assessable for a specified task.

The intersection of several categories deserves the same assessment.
Joint grammatical, expressive, and interactional membership may be an
economical grouping for some questions. Its derivability from the separate
memberships neither establishes nor defeats that possibility. Conversely,
an advantage over an artificially restricted additive model wouldn't
establish the intersection as a privileged category for every inquiry.
The comparison should identify what the competing representations already
contain.

Historical explanation poses a different question. If a documented
recruitment process produced a present property, controlling for that
property can remove the association through which the process is recognized.
The disappearance of an additional predictive effect then leaves the
historical explanation open. A test of prediction from limited present-day
evidence and a test of the proposed history have different evidential
requirements. Section~\ref{sec:history} develops the distinction through
\mention{gee}.

Processing and social explanations require similar care about which
variables are compared. Repeated use could help produce fluent retrieval;
controlling all frequency-related variation would answer a question about
effects beyond frequency. It wouldn't by itself decide whether repetition
helped establish the routine. A social classification may concern
audience-sensitive selection, making the relation to audience part of the
target rather than a nuisance to eliminate.

These explanations also differ in strength. A recurring profile can be
stable while its causal ordering remains unknown. An independently
supported historical dependence can be real without demonstrating that
the whole profile recurs across an inventory. Maintenance requires evidence
that a process preserves a profile, and corrective control requires a
response to disturbance. I reserve \term{homeostatic} for that stronger
control claim, following \textcite{reynolds2026notEveryStableCluster}, a
narrower use than Boyd's \citep[143]{boyd1999}. No such control claim is
established here.

The role of inquiry is consequently limited and substantive. It selects
the bearer, the contrast of interest, the scope, and the standard of success.
Observed distributions and supported explanations constrain which groupings
meet that standard. This permits several justified classifications without
making category standing a consequence of an investigator's preference.
